\part{Requirements, architecture and end-to-end data flow}

\section{Requirements engineering}

\begin{qa}{97}{What are the functional requirements of CARS-SF?}
\begin{enumerate}[label=F\arabic*,leftmargin=2em]
\item Infer the operating context from clock, calendar, alerts and admin input.
\item Map traffic to a service slice (by address, application or prediction).
\item Estimate per-slice demand from measurements.
\item Compute a feasible allocation with floors.
\item Enforce it in both directions of the bottleneck.
\item Monitor per-slice SLAs.
\item Detect, diagnose and recover from backhaul and core faults.
\item Fail back safely.
\item Log every decision with its inputs.
\item Produce measured datasets for learning.
\end{enumerate}
\end{qa}

\begin{qa}{98}{What are its non-functional requirements?}
\begin{enumerate}[label=N\arabic*,leftmargin=2em]
\item Detection in under 1\,s; critical outage under 1\,s.
\item Re-plan within one control interval (5\,s).
\item Never form an L2 loop.
\item Learned components cannot bypass floors.
\item Negligible control-plane overhead.
\item Reproducible runs (seeded, versioned).
\item Swappable stages.
\item Sudo-free evaluability.
\end{enumerate}
\end{qa}

\begin{qa}{99}{Which requirements are measurable, and which are qualitative?}
Measurable: N1 (MTTD/MTTR, outage), N2 (controller timestamps), N3 (by construction plus tests), N5 (probe and tc
command rates), N6 (re-running a seed gives the same traffic schedule; unit test). Qualitative: N7 (demonstrated by
three QoS backends and seven schedulers), F2--F10 (demonstrated, not scored).
\end{qa}

\begin{qa}{100}{How did you prioritise requirements?}
By harm. Safety first (no loops, floors respected, no fail-back onto a bad link), then critical SLA, then efficiency
(surplus), then extensibility (ML, 5G).
\end{qa}

\begin{qa}{101}{Which requirements are safety-critical?}
N3 (a loop would cause a broadcast storm and take down the village network), N4, the floors, and safe fail-back
(OAM-gated). All are enforced by construction and covered by unit tests (\code{tests/test_self_healing_engine.py}).
\end{qa}

\begin{qa}{102}{System requirement vs implementation detail?}
``Critical outage under 1\,s after a backhaul cut'' is a requirement. ``50\,ms link-watcher polling with an event that
wakes the loop'' is one implementation meeting it.
\end{qa}

\begin{qa}{103}{What is the expected monitoring interval?}
Telemetry and SLA monitoring every 1\,s; RTT probes once per second per service, in a background thread; link state
every 50\,ms; re-planning every 5\,s.
\end{qa}

\begin{qa}{104}{What maximum detection and recovery time is acceptable for critical services?}
Target: detection under 1\,s and service outage under 1\,s, roughly the tolerance of interactive voice and video
(beyond that, calls drop or freeze). Achieved: MTTD about 0\,s with the link watcher; outage \NOutCutFullMean\,s
(Linux bridges) and \NOvsOutageAfter\,s on real OVS.
\honest{The 1\,s budget is a design choice, not derived from a clinical standard.}
\end{qa}

\begin{qa}{105}{What availability target would you set, and how would you measure it?}
For critical slices, 99.9\% of 1-minute intervals SLA-compliant, measured by \code{SliceMonitor} over long runs.
\honest{Not measured: runs are 40--45\,s, so we report outage per failure event instead.}
\end{qa}

\begin{qa}{106}{How much control-plane overhead is acceptable?}
Below 1\% of bottleneck capacity and a few percent of one CPU core. Measured costs:
\begin{itemize}
\item four ICMP probes per second (negligible bandwidth);
\item 8 \code{tc class change} commands per re-plan;
\item a controller step of about \NCtrlStepUs\,$\mu$s; a PPO re-plan of about \NPpoReplanMs\,ms.
\end{itemize}
Endpoint captures are \emph{evaluation} overhead, not a control-plane requirement.
\end{qa}

\begin{qa}{107}{How should the system behave when requirements conflict?}
By a fixed precedence:
\begin{enumerate}
\item safety (loop-freedom, floors);
\item healing constraints (no core slice above backup capacity);
\item context priorities;
\item demand and efficiency.
\end{enumerate}
For context, admin override beats alerts, which beat the calendar.
\end{qa}

\begin{qa}{108}{Which requirements are satisfied, partially satisfied, or unverified?}
\begin{itemize}
\item \textbf{Satisfied (L3):} F1--F10 except as below; N1--N4; N6; N8.
\item \textbf{Partial:} F7 (no generic post-action verify-and-rollback); F9 (no correlation IDs across artefacts);
N5 (not stress-tested).
\item \textbf{Unverified:} availability over time; behaviour at scale (hundreds of nodes); behaviour under real radio.
\end{itemize}
\end{qa}

\begin{qa}{109}{How are requirements mapped to modules and test cases?}
\begin{center}\scriptsize
\begin{tabular}{@{}lll@{}}\toprule
Req. & Module & Evidence \\\midrule
F1 & \code{context/engine.py} & \code{test_context.py}, timeline study \\
F3 & \code{slicing/demand_estimator.py} & \code{test_slicing.py} (probe-up, Holt) \\
F4 & \code{slicing/allocators.py} & \code{test_slicing.py} (floors, capacity) \\
F5 & \code{slicing/qos_enforcer.py} & \code{test_qos.py} (exact tc commands); all runs \\
F7, F8, N3 & \code{self_healing/engine.py} & \code{test_self_healing_engine.py}; resilience study; real OVS \\
F10 & \code{features/} & \code{test_features.py} (OWD, loss on synthetic pcaps) \\
N1 & engine + fast path & \code{test_link_event_fast_path...}; OVS MTTR \NOvsMttrAfter\,s \\
\bottomrule\end{tabular}\end{center}
\end{qa}

\begin{qa}{110}{What would be the acceptance criteria for declaring CARS-SF complete?}
\begin{itemize}
\item All unit and integration tests green.
\item The full matrix reproducible from a clean checkout.
\item Every hypothesis evaluated, positive or negative.
\item Each claim in the thesis carrying a maturity level.
\item One sudo run on real OVS for each failure type.
\item For a deployment claim: an L4 field trial.
\end{itemize}
\end{qa}

\section{High-level architecture}

\FigArch

\begin{qa}{111}{Draw the complete architecture from traffic generation to recovery verification.}
Fig.~\ref{fig:arch}. The tick order is in Fig.~\ref{fig:tick}, and ``recovery verification'' is the state machine's
escalation and fail-back guards (Fig.~\ref{fig:heal}).
\end{qa}

\FigTick

\begin{qa}{112}{Why did you choose a modular architecture?}
\why{Research needs to swap one mechanism at a time (allocator, policy, QoS backend) and compare.}
\choice{Protocol-typed stages (\code{core/interfaces.py}) composed by name from \code{schedulers.yaml}.}
\alt{A monolithic controller (faster to write, impossible to ablate cleanly), or microservices (operational overhead).}
\trade{Ablations and baselines are config changes; the price is indirection and interfaces that are only checked by
duck-typing.}
\ev{Seven schedulers and three QoS backends (tc, OVS, UPF) run on the same controller.}{\Lthree}
\end{qa}

\begin{qa}{113}{Responsibilities of each module.}
\begin{itemize}
\item \textbf{Context}: infers the context id.
\item \textbf{Traffic generator}: model-driven, seeded flows (apparatus).
\item \textbf{Topology}: builds the network and owns the path primitives.
\item \textbf{Telemetry}: per-tick deltas, class stats, probes, link events.
\item \textbf{Features}: packet matching and windows (post-run; shape features also live).
\item \textbf{Dataset}: builds and validates the dataset.
\item \textbf{Classifier}: SliceNet, offline and in-loop.
\item \textbf{Allocation}: rates.
\item \textbf{QoS}: tc / OVS / UPF commands.
\item \textbf{Monitor}: per-slice SLA records.
\item \textbf{Detector}: failure and prediction events.
\item \textbf{Planner / healer}: path actions and constraints.
\item \textbf{Verifier}: escalation, hold-down and OAM guards inside the healer (no separate module).
\end{itemize}
\end{qa}

\begin{qa}{114}{Which modules are on the live decision path?}
Context, demand estimator, priority engine (plus an optional CEM or PPO policy), allocator, QoS enforcer, telemetry,
SLA monitor, detector, predictor and healer. The in-loop classifier is on the path only when
\code{classification=slicenet}.
\end{qa}

\begin{qa}{115}{Which modules operate offline?}
Packet matching, feature extraction and dataset building (after each run), the validator, SliceNet training, RL
training and calibration, aggregation and reporting, and the RSNSD generator.
\end{qa}

\begin{qa}{116}{Which modules are experimental alternatives rather than the active production flow?}
CEM and PPO policies, the no-surplus, proportional, static, 3GPP and no-slicing allocators, the OVS QoS backend, the
oracle demand mode, the legacy \code{RecoveryPlanner}, and the legacy \code{MininetSims/} and \code{RSNSDv2/}
(excluded from tests). The default (production) composition is the \texttt{proposed} scheduler with the tc backend
and host classification.
\end{qa}

\begin{qa}{117}{What is the canonical source of truth for the architecture?}
The code plus the registries. Prose: \code{ARCHITECTURE.md}, whose §12 maps every original architecture module to
code with its status. \code{architecture.md} is the original vision, kept with a status banner.
\end{qa}

\begin{qa}{118}{How do you prevent the architecture document from describing capabilities that aren't integrated?}
\begin{itemize}
\item The status table states maturity (load balancing is marked \emph{not implemented}).
\item Every claim in the results is generated from data (\code{report.py}, \code{defense_numbers.py}).
\item Integration tests assert the integrated behaviours (context switching, predictive failover, fault localisation,
in-loop classification).
\end{itemize}
\honest{The original document did over-claim (``QoS Enforcer: prototype''). The verification pass that started
this work found the Mininet path had never run.}
\end{qa}

\begin{qa}{119}{Which module owns the current state of the network?}
No single owner. Telemetry owns the observed state (snapshot per tick). The controller owns the intended allocation
(\code{controller.allocation}). The healer owns the path state (\textsf{NORMAL} / \textsf{ON\_BACKUP} /
\textsf{DEGRADED}). The kernel owns the actual state.
\honest{Intended and actual state are not reconciled by reading back tc (Q139).}
\end{qa}

\begin{qa}{120}{Which module owns the service requirements and SLA definitions?}
\code{registries/slices.yaml} (slice SLAs, floors, class ids) and \code{registries/applications.yaml} (per-application
SLAs and models). They are read by the monitor, the extractor (SLA labels), the allocator (floors) and QoS (classes).
\end{qa}

\begin{qa}{121}{How are modules coupled?}
\begin{itemize}
\item Direct in-process calls inside the orchestrator loop.
\item YAML configuration.
\item Two threads (probe worker, link watcher) with locks, plus a \code{threading.Event} that wakes the loop.
\item Subprocesses (generators, sinks, tcpdump).
\item Shell commands to the network (\code{tc}, \code{ip}, \code{ovs-*}).
\item Files for post-run artefacts.
\end{itemize}
\end{qa}

\begin{qa}{122}{Why is that communication mechanism appropriate?}
One host, one experiment, millisecond control timescales: in-process calls are the simplest and fastest, and keep
runs deterministic and debuggable. A broker would add latency and failure modes without a distribution need.
\end{qa}

\begin{qa}{123}{What happens when one module is slow, unavailable, or returns invalid data?}
\begin{itemize}
\item Probes run in a background thread, so a lost probe never blocks the loop; it reports \texttt{None}.
\item Telemetry readers catch errors and return empty or zero values.
\item Ticks are wall-clock aligned, so a late tick does not shift later ones.
\item A missing RL model falls back to rule-based, and this is recorded (\code{policy_fallback_to_rule}).
\item A missing SliceNet model disables in-loop classification.
\end{itemize}
\honest{tc command failures are not detected (exit codes are not checked), and missing telemetry can look like zero
traffic (Q279).}
\end{qa}

\begin{qa}{124}{Which interfaces have explicit input and output schemas?}
The stage protocols in \code{core/interfaces.py} (types only). The dataset schema (\code{SCHEMA_COLUMNS} and
\code{registries/features.yaml}, validated by \code{DatasetValidator}). The scheduler registry is validated by name.
\honest{The telemetry snapshot is a convention-based dict and the YAML registries have no JSON-schema validation.}
\end{qa}

\begin{qa}{125}{How are units, timestamps, identifiers and missing values represented?}
\begin{itemize}
\item Units are in the names (\code{_mbps}, \code{_ms}, \code{_pct}, \code{_bytes}).
\item Timestamps are Unix epoch seconds (float) from one host clock.
\item Identifiers: \code{run_id}, \code{application_id}, \code{ue_host}, \code{flow_id}, \code{window_index}.
\item Missing values: \texttt{None} in telemetry (a lost probe), \texttt{NaN} in datasets (no delivered packets).
\end{itemize}
\end{qa}

\begin{qa}{126}{How are configurations validated before execution?}
Scheduler names are checked against the registry; CLI options are restricted by argparse \code{choices}; scenario
and timeline ids raise on unknown values; the allocator and QoS raise on unknown classes.
\honest{Value ranges in YAML (negative demand, floors above capacity) are not validated up front.}
\end{qa}

\begin{qa}{127}{How do you prevent inconsistent definitions of service priority across modules?}
One source: \code{scenarios.yaml} priority weights, read by \code{slicing/context_engine.py}, which overlays the
registry on its legacy table. HTB prio and class ids come from \code{slices.yaml}.
\honest{\code{SERVICE_CLASS_IDS} is duplicated in the QoS enforcer and the telemetry collector, with the same values;
it should be read from \code{slices.yaml}.}
\end{qa}

\begin{qa}{128}{How are configuration versions recorded with experiment results?}
\code{run_meta.json} stores all parameters plus the git commit, and \code{summary.json} stores the scheduler spec and
all options.
\honest{Uncommitted registry edits are not captured by the hash. The matrix was always run from committed code.}
\end{qa}

\begin{qa}{129}{What are the major architectural failure points?}
\begin{itemize}
\item The single orchestrator process (no high availability).
\item Host CPU contention distorting timing (observed and fixed).
\item Silent tc failures.
\item Telemetry stalls that look like zero traffic.
\item The link watcher depending on carrier propagation (which differs between OVS and Linux bridges).
\item Disk for captures (about 100\,MB per run).
\end{itemize}
\end{qa}

\begin{qa}{130}{What would you redesign first to support hundreds of nodes?}
\begin{itemize}
\item Split into per-site agents (local enforcement plus fail-safe defaults) and a regional controller.
\item Replace the polling of \code{/sys} and \code{tc} with streaming telemetry (sFlow, IPFIX, gNMI).
\item Keep state in a store with leases.
\item Use OpenFlow or P4 path control instead of admin-state toggling.
\item Make allocation a per-link optimisation.
\end{itemize}
\end{qa}

\section{End-to-end execution}

\begin{qa}{131}{Walk one application flow from creation to final QoS enforcement.}
\begin{enumerate}
\item The \code{TrafficOrchestrator} plans the telemedicine app on \code{h_health} with demand-scaled parameters and a
fixed seed.
\item Before traffic starts, the QoS setup installs HTB classes and a host filter \code{src 10.0.1.10 → 1:10} on
\code{s_rural-up} (and \code{dst} on \code{s_edge-down}).
\item The generator sends VBR UDP; the packets hit the u32 filter and the fq\_codel leaf of class 1:10 at the current
rate and ceiling.
\item Telemetry counts bytes; the controller re-plans every 5\,s and issues \code{tc class change}.
\item After the run, captures on \code{h_health} and \code{srv_edge} are matched, giving one-way delay and loss per
5\,s window, labelled \code{healthcare.telemedicine}.
\end{enumerate}
\end{qa}

\begin{qa}{132}{Walk through a medical-emergency scenario from context change to verified recovery.}
Timeline \code{medical_night}: tick 6 raises \code{medical_emergency}.
\begin{enumerate}
\item The context engine returns the new context with \code{changed=True}.
\item The traffic orchestrator re-spawns generators with the emergency demand (healthcare 95\,Mb/s).
\item The controller re-plans immediately (reason: context change) with weights 15/1/1/1 from the prior demand,
because the estimator restarts with a warm-up.
\item The QoS classes change.
\item The SLA monitor records the healthcare slice every tick; latency spikes trigger $\times$2 boosts for 6 ticks.
\end{enumerate}
\honest{``Verified recovery'' here means the SLA monitor shows compliance afterwards. There is no explicit ``SLA
restored'' event that ends the episode.}
\end{qa}

\begin{qa}{133}{Which component first observes the emergency?}
The context engine, through an alert event in the timeline (in production: an alert feed or operator). The network
does not ``observe'' an emergency; it observes its consequences (demand, RTT).
\end{qa}

\begin{qa}{134}{How does the emergency change service priorities?}
The priority matrix switches to the context's weights (healthcare 15 vs 1), and the allocator gives healthcare its
weighted share up to demand. The floors stay unchanged.
\end{qa}

\begin{qa}{135}{How is traffic associated with a service class?}
\begin{itemize}
\item \textbf{host} (default): UE address $\to$ service.
\item \textbf{app}: registry server and port (oracle; also used for 3GPP slice types).
\item \textbf{slicenet}: per-flow prediction for unknown or mixed UEs.
\end{itemize}
The dataset labels come from spawn metadata (ground truth).
\end{qa}

\begin{qa}{136}{How does that classification reach the allocation engine?}
It doesn't need to. Classification decides \emph{which HTB class} packets enter; the allocator decides \emph{each
class's rate}. They meet in the kernel. Demand per class is measured per class.
\end{qa}

\begin{qa}{137}{How does the allocation engine's decision become network configuration?}
\code{SliceController.step} $\rightarrow$ \code{QoSEnforcer.apply_allocation}, which issues
\code{tc class change … rate X ceil Y} on \code{s_rural-up} and \code{s_edge-down} (ceiling = capacity, or the
allocation for core slices while on backup). On OVS: \code{ovs-vsctl set queue …}. On the 5G UPF: the same tc commands
via \code{docker exec}.
\end{qa}

\begin{qa}{138}{How does monitoring determine whether the configuration worked?}
Indirectly, through outcomes: per-class throughput (counters), drops (HTB stats), probe RTT, and SLA compliance per
slice per tick. Post-run, packet-matched KPIs per window.
\end{qa}

\begin{qa}{139}{What if the allocation decision is correct but the network configuration fails?}
\honest{Today it goes undetected at configuration time. The controller assumes success, and only the SLA monitor would
show the violation later.}
Fix (small): after \code{apply_allocation}, read back \code{tc -s class show} (\code{get_queue_stats} exists),
compare rate and ceiling, retry once, and raise a \code{config_failure} event to the healer or operator.
\end{qa}

\begin{qa}{140}{What if the configuration succeeds but the SLA remains violated?}
The detector raises \code{latency_spike} (RTT above 200\,ms) or \code{congestion} (more than 50 drops per tick). The
healer boosts the affected slice's priority $\times$2 for 6 ticks (with a cooldown), and the controller re-plans.
The predictor may also flag a backhaul trend.
\honest{There is no escalation if repeated boosts fail (Q508).}
\end{qa}

\begin{qa}{141}{How do you prevent a module from reporting success when a downstream module has failed?}
\begin{itemize}
\item Runs are judged on measured data-plane outcomes, never on module return codes.
\item Integration tests assert data-plane facts (a PCAP dataset exists, the outage was measured, actions were taken).
\item A run whose dataset falls back to telemetry is marked \code{degraded} by the matrix runner (none of the 340).
\end{itemize}
\end{qa}

\begin{qa}{142}{What evidence shows the complete path works, not just each component?}
\begin{itemize}
\item The 7 integration tests (timeline, failover, prediction, core localisation, baselines, in-loop SliceNet).
\item 340 matrix runs, every one producing packet-matched data.
\item Real-OVS runs in a privileged container.
\item The 5G-core runs.
\end{itemize}
All are L3 in the named environment.
\end{qa}

\begin{qa}{143}{Which exact command or entry point starts the experiment?}
\begin{itemize}
\item Real OVS: \code{sudo .venv/bin/python run_experiment.py --scenario fiber_cut --duration 45 --apply-qos}.
\item Sudo-free: \code{python -m tests.integration.harness fiber_cut 45 apply_qos=true}.
\item Studies: \code{experiments/run_matrix.py}.
\item 5G: \code{deployment/5g/run_5g.py}.
\end{itemize}
\end{qa}

\begin{qa}{144}{Which processes, threads, namespaces, interfaces and background services are involved?}
\begin{itemize}
\item The orchestrator process with a probe thread and a link-watcher thread.
\item A Mininet shell per node.
\item 12 network namespaces (8 UEs and 4 servers) plus the root namespace for switches.
\item About 24 veth pairs; about 40 generator and sink processes; 12 tcpdump processes.
\item On OVS: \code{ovsdb-server} and \code{ovs-vswitchd}.
\end{itemize}
The harness wraps all of this in one user+net+mount namespace.
\end{qa}

\begin{qa}{145}{What artefacts are generated after a run?}
\code{results/runs/<run_id>/}: \code{summary.json}, \code{dataset.csv}, \code{sequences.npz}, \code{feedback.csv},
\code{slice_monitor.csv}, \code{validation_report.txt}, \code{run_meta.json}. Also captures in
\code{generated/pcap/<run_id>/} and generator logs in \code{generated/logs/<run_id>/}.
\end{qa}

\begin{qa}{146}{How do you trace one decision across logs, telemetry, configuration and metrics?}
By tick and timestamp: \code{controller_history} (tick, reason, demand estimate, priorities, allocation, applied)
$\rightarrow$ \code{feedback.csv} (per tick: allocation, throughput, RTT) $\rightarrow$ dataset windows
(\code{window_index}, flow times) $\rightarrow$ healing actions (timestamps, trigger).
\honest{There is no shared decision ID; correlation is by time.}
\end{qa}

\begin{qa}{147}{Data pipeline that produces a dataset vs runtime pipeline that controls a network?}
The runtime loop acts every second on aggregate telemetry. The data pipeline reconstructs per-flow, per-window facts
after the run from packet captures (ground-truth labels, one-way delay). The first is causal and online; the second is
retrospective and exact.
\end{qa}

\begin{qa}{148}{Which modules are shared by both pipelines?}
Topology, traffic models, telemetry, the registries, and \code{features/shape.py} (shape features and sequences, used
both for training and for the in-loop classifier).
\end{qa}

\begin{qa}{149}{Where can their behaviour diverge?}
\begin{itemize}
\item The in-loop classifier sees one vantage point (the UE capture) and the cumulative flow history; training used
5\,s windows from the sender side of matched packets.
\item Rate features are normalised by the observed span (at least one window) to reduce the mismatch.
\item Live telemetry is aggregate; dataset QoS is per flow.
\end{itemize}
\end{qa}

\begin{qa}{150}{How do you ensure the runtime uses the same feature definitions and taxonomy as offline training?}
The same functions (\code{shape_features}, \code{packet_sequence}); the trained model stores its feature column order
and normalisation; the four-class taxonomy comes from one registry. The in-loop classifier picks whichever model scored
higher on held-out data.
\honest{There is no automatic train/serve skew check.}
\end{qa}

\section{Architectural alternatives and trade-offs}

\begin{qa}{151}{Why not implement everything in one Python process?}
It nearly is: one orchestrator process with threads. Generators and tcpdump are separate processes because they must
run \emph{inside} host namespaces and must not block the loop.
\end{qa}

\begin{qa}{152}{Why not use a message broker between every module?}
\alt{Kafka or MQTT between modules.}
\trade{It would decouple and distribute, but add milliseconds, ordering issues and a broker as a single point of
failure, with no need on one host. It becomes justified with per-site agents (Q130).}
\end{qa}

\begin{qa}{153}{Why not use an OpenFlow SDN controller as the central component?}
\why{The problem is shaping and scheduling at a bottleneck and path protection. OpenFlow can steer packets to queues
but does not schedule them.}
\choice{The kernel qdiscs do the scheduling. The topology uses \code{OVSBridge} in standalone mode (no controller);
the OVS backend uses OpenFlow \code{set_queue}.}
\alt{Ryu or ONOS with meters and queues. It would add path computation for multipath healing.}
\trade{Simpler and fewer moving parts; the cost is no flow-level path control (hence no load balancing).}
\end{qa}

\begin{qa}{154}{Why not use a rule engine instead of an ML decision layer?}
We essentially \emph{do}: the default decision layer is rule-based (context weights, floors, water-filling,
threshold detectors). ML is used only where rules lack information (traffic identity for a mixed UE) and as an
evaluated alternative (RL). It is not justified by results.
\end{qa}

\begin{qa}{155}{Why not use a centralised optimiser with a global view?}
For a single link with a separable concave objective, weighted water-filling with caps is essentially what an LP or
convex solver returns (Q432). A global optimiser pays off with multiple links and paths. The global view here is
just one link.
\end{qa}

\begin{qa}{156}{Advantages and disadvantages of centralised, distributed and hierarchical control?}
\begin{itemize}
\item Centralised: simple, consistent, but a single point of failure and a partition risk.
\item Distributed: resilient, local, but needs coordination for consistent shares.
\item Hierarchical: local enforcement and fast failover at the site, regional policy above.
\end{itemize}
CARS-SF is centralised per site, and that is the right granularity for one village. Regional deployment should be
hierarchical.
\end{qa}

\begin{qa}{157}{How would you handle a partition that isolates the controller from some switches?}
Switches keep their last HTB configuration (the kernel retains it), so enforcement continues. Healing on the cut-off
side stops.
\honest{No local fail-safe agent exists today. Add one with a lease: if the controller is silent for N seconds, revert
to the context-independent safe profile and run local failover.}
\end{qa}

\begin{qa}{158}{What should keep operating locally if the central controller is unavailable?}
Floors and the last allocation (already persistent in the kernel); local link-failure switchover (protection
switching should be local by design); loop-safety interlocks.
\end{qa}

\begin{qa}{159}{How would you add a new service class without modifying every module?}
Registry entries in \code{slices.yaml}, \code{applications.yaml} and \code{scenarios.yaml} are data-driven for QoS
classes, monitoring and traffic.
\honest{Several modules hard-code the 4-tuple \code{SERVICE_CLASSES} (allocators, env, analysis). Generalising that
tuple to read the registry is required.}
\end{qa}

\begin{qa}{160}{How would you add a new telemetry backend without rewriting allocation logic?}
Implement \code{TelemetrySource} so that it returns the same snapshot keys (\code{service_throughput_mbps},
\code{tc_class_stats}, \code{service_rtt_ms}, \code{link_events}…). Allocation consumes only those through the demand
estimator.
\end{qa}

\begin{qa}{161}{Which abstractions are essential, and which add unnecessary complexity?}
Essential: the allocator, policy, QoS backend and telemetry snapshot. Questionable: the legacy \code{RecoveryPlanner}
(superseded but kept for old tests) and the two context classes (\code{ContextManager} is a thin logger).
\end{qa}

\begin{qa}{162}{How would you separate policy, decision-making and enforcement more cleanly?}
Policy = registries plus a policy object (weights, floors, SLAs) with a version. Decision = a pure function
\code{(state, policy) → intent}. Enforcement = backend adapters that \emph{acknowledge} and read back. Today the
healer mixes decision and enforcement (it calls topology primitives directly).
\end{qa}

\begin{qa}{163}{What is the smallest set of interfaces that preserves modularity?}
Four: \code{TelemetrySource.poll}, \code{Allocator.allocate}, \code{QoSBackend.apply_allocation} and
\code{Healer.handle/constraints}. Context and policy can be functions.
\end{qa}

\begin{qa}{164}{Which design decisions were driven by research needs rather than production requirements?}
\begin{itemize}
\item endpoint packet capture (evaluation ground truth);
\item seeded generators;
\item the oracle demand mode;
\item seven baselines;
\item the sudo-free namespace harness;
\item per-run artefact directories;
\item the predictive detector's OAM shortcut (reading the emulated impairment).
\end{itemize}
\end{qa}
